\section*{ID8648: Relation: 81/80≈21.5}
\paragraph{Metadata.}
PiID: 4197484423608 | RTEID: RTE:P33582.1045:T5.0675 | UUID: eefed4c3-03a0-5519-a2ba-cfc13f8a13d9

\paragraph{Canonical Statement.}
Musically, the combination of \$16/15\$ and 10/9 is significant. It shows how an 18.5185\% frequency increase corresponds to a minor-third interval in one tuning system. Interestingly, equal temperament's minor third is \$2\textasciicircum{}\{3/12\}\textbackslash{}approx1.1892\$ (300 cents), so the Pythagorean value 1.185185 is about 6 cents flat of equal temperamenten.wikipedia.org. This subtle difference (caused by the syntonic comma, \$81/80\textbackslash{}approx21.5\$ centsen.wikipedia.org) highlights the beating or "instability" that early theorists noticed in the \$32/27\$ intervalmedieval.org. In medieval music, 32:27 was called the semiditone and considered the "mildest level of dissonance" for a thirdmedieval.org - being more tense than the pure 6:5 minor third but still usable. Thus, 1.185... links directly to the language of musical ratios: it encapsulate

\paragraph{Canonical Equation.}
\[
81/80\approx21.5
\]

\paragraph{Dictionary.}
Relation: 81/80≈21.5 — 81/80≈21.5

\paragraph{Encyclopedia.}
Relation: 81/80≈21.5 is a scaling / order-of-magnitude relation, written 81/80≈21.5. It expresses a scaling or proportionality, capturing how the quantity grows with its parameters. Within the Trawin Zero Point registry it serves as one element of the framework's mathematical narrative.
